\section{Global finite-period degeneration of the original fixed conductor spaces}
Fix every original quartet and nonzero unit allowed in ANP1, including
its complete four-factor period matrix and all original coefficients.
The exact original norms, translations and Gamma frames below are those
of FC20--24 and WCF1--6 \cite{WCF,Cumulative,PCLOriginal}.
No five-orbit hypothesis is imposed in this section. In particular,
the symbol may vanish identically at the selected period.

Write $F_j(z)=f_j(0,z)$ for the four original functions and
\[
 E_A(\xi,z)=\prod_{j=1}^4 f_j(\xi,z),\qquad
 \mu_m(z)=\partial_\xi^mE_A(0,z),\qquad
 \mu_0(z)=\prod_{j=1}^4F_j(z).
 \tag{GFC1}
\]
ANP2 proves that each $F_j$ is nonconstant, nonzero entire and has a
nonzero finite zero. Thus $\mu_0$ is a nonzero entire function with a
nonzero finite zero $z_0$. The product cannot vanish identically:
the zeros of each nonzero entire factor are isolated, and the union
of four such zero sets cannot contain a disk.
Choose any such $z_0$, not necessarily a simple or isolated-factor zero.
For the complete four-factor multiplicities and coefficients set
\[
 \ell_j=\min\{m\ge0:F_j^{(m)}(z_0)\ne0\},\quad
 \nu=\sum_{j=1}^4\ell_j\ge1,\quad
 c=\prod_{j=1}^4\frac{F_j^{(\ell_j)}(z_0)}{\ell_j!}\ne0.
 \tag{GFC2}
\]
All four integers are finite because their factors are nonzero entire.
With $\tau=z-z_0$, multiplication of the full Taylor series gives
\[
 \mu_0(z)=c\tau^\nu+O(\tau^{\nu+1}),\qquad
 c=\frac{\mu_0^{(\nu)}(z_0)}{\nu!}.
 \tag{GFC3}
\]
An isolated-zero disk can be chosen small enough that $z\ne0$ and
$\mu_0(z)\ne0$ for $0<|\tau|<\epsilon$.

\subsection{The complete fixed-space matrix and its entire inverse-exterior consequences}
For the unchanged cutoff $D\ge0$ put $N=D+1$ and retain the original
$s\in\{1,a_{\rm amp}\}$, $M_s=(2\pi)^{s/2}$,
$\rho_n=\sqrt{n!/(s/2)_n}$, and centers
$c_0=a_{\rm amp}/2$, $c_1=a_{\rm amp}/2-4$.
The two fixed old spaces are
\[
 (\mathcal P_D,\|\cdot\|_{s,c_0})
       \longrightarrow(\mathcal P_D,\|\cdot\|_{s,c_1}),\qquad
 \|P\|_{s,c}^2=\int_{\mathbb R}|P(c+iy)|^2
 \frac{(2\pi)^{s/2}2^{s/2}}{4\pi\Gamma(s/2)}
       |\Gamma(s/4+iy/2)|^2\,dy.
 \tag{GFC4}
\]
The original orthonormal frames are
$\phi_{n,s,c}=p_{n,s}((S-c)/i)/\sqrt{M_sn!(s/2)_n}$,
with the full WCF2 recurrence for $p_{n,s}$.
Define
\[
 \omega(t)=-i\arctan t,\qquad
 g(t,z)=e^{-4\omega(t)}E_A(\omega(t),z)
       =\sum_{m\ge0}g_m(z)t^m,
 \quad
 R_D=\operatorname{diag}(\rho_0,\ldots,\rho_D).
 \tag{GFC5}
\]
All $g_m(z)$ are entire: each is a finite combination of original
moments, and every original coefficient is entire in $z$ by PCL2--10.
This is the fixed-space symbol, without a division by a symbol-order
power. It has $g_0(z)=\mu_0(z)$ at every period.
The exact matrix of the original translation conductor on GFC4 is
\[
 B_{D,s}(z)=R_D^{-1}T_D(g(\cdot,z))R_D,
 \qquad
 B_{ij}(z)=
 \begin{cases}
 g_{j-i}(z)\rho_j/\rho_i,&0\le i\le j\le D,\\
 0,&i>j.
 \end{cases}
 \tag{GFC6}
\]
The original generating identity WCF3--6 proves these entries without
requiring a nonzero symbol. The common mass cancels in the two
orthonormal-frame entries, while the norms in GFC4 retain it.
In particular every diagonal entry is the unchanged $\mu_0(z)$, so
\[
 \det B_{D,s}(z)=\mu_0(z)^N,\qquad
 B_{D,s}(z)\text{ is invertible exactly when }\mu_0(z)\ne0.
 \tag{GFC7}
\]
The coordinate matrix has characteristic polynomial
$(\lambda-\mu_0(z))^N$. This is a statement in the two specified
orthonormal coordinate frames; no identification of their physical
centers is required to use the following norm consequence.

For every $1\le r\le N$ and every $\mu_0(z)\ne0$,
\[
 \boxed{\qquad
 \|\bigwedge^r B_{D,s}(z)^{-1}\|
                \ge |\mu_0(z)|^{-r}.
 \qquad}\tag{GFC8}
\]
To prove this directly, the inverse of the upper triangular matrix
GFC6 is upper triangular with diagonal $\mu_0^{-1}$.
In the orthonormal exterior power frames, each diagonal matrix entry
of its $r$th exterior power is the corresponding principal minor,
equal to $\mu_0^{-r}$. The norm of an operator is at least the
absolute value of any one of its matrix entries in orthonormal
frames: apply it to the unit source basis vector and project onto
the unit target basis vector. This proves GFC8 with all original
Gamma weights and centers retained.
For the top exterior power the space is one-dimensional, giving
the stronger exact identity and pole:
\[
 \|\bigwedge^N B_{D,s}(z)^{-1}\|=|\mu_0(z)|^{-N},\qquad
 \lim_{z\to z_0}|z-z_0|^{\nu N}
       \|\bigwedge^N B_{D,s}(z)^{-1}\|=|c|^{-N}>0.
 \tag{GFC9}
\]
GFC8 also gives, at every exterior rank,
$\liminf_{z\to z_0}|z-z_0|^{\nu r}
\|\bigwedge^rB_{D,s}(z)^{-1}\|\ge|c|^{-r}$.

There is a complete finite test for the exact pole and coefficient at
every rank, even when the symbol vanishes identically at $z_0$.
For $0\le k\le N$ let
\[
 m_k=\min\left\{0\le m\le\nu k:
     \left.\frac1{m!}\partial_z^m\bigwedge^kB_{D,s}(z)
          \right|_{z=z_0}\ne0\right\},\qquad
 L_k=\left.\frac1{m_k!}\partial_z^{m_k}
                  \bigwedge^kB_{D,s}(z)\right|_{z=z_0}.
 \tag{GFC10}
\]
For $k=0$ these definitions mean $m_0=0,L_0=1$.
For every positive $k$ the finite test terminates: a principal minor
of $B$ equals $\mu_0^k$, whose order is exactly $\nu k$.
Thus the indicated exterior matrix has a nonzero coefficient no
later than that order. Every derivative in GFC10 is a finite
determinant expansion in the explicit original entries GFC6 and
their derivatives, so this test introduces no free matrix or norm.

For $1\le r\le N$ put $k=N-r$. The singular values of an
invertible $N$ by $N$ matrix give the exact identity
\[
 \|\bigwedge^rB^{-1}\|
       =\frac{\|\bigwedge^{N-r}B\|}{|\det B|}.
 \tag{GFC11}
\]
Indeed if the forward singular values are
$\sigma_1\ge\cdots\ge\sigma_N>0$, its numerator is
$\sigma_1\cdots\sigma_{N-r}$ and its denominator is
$\sigma_1\cdots\sigma_N$. Their quotient is the product of
the $r$ largest inverse singular values, proving the identity.
Since $\bigwedge^kB=\tau^{m_k}(L_k+O(\tau))$, continuity
of the operator norm now proves the full exact pole data
\[
 \boxed{\quad
 p_r=\nu N-m_{N-r}\ge\nu r,\qquad
 \lim_{z\to z_0}|z-z_0|^{p_r}
       \|\bigwedge^rB_{D,s}(z)^{-1}\|
       =\frac{\|L_{N-r}\|}{|c|^N}>0.
 \quad}\tag{GFC12}
\]
This gives integer pole exponents and their exact original-metric
constants at all ranks; GFC9 is the case $r=N$.

\subsection{Every possible symbol order at the selected zero}
The nine-frequency test and order bound $32$ are the earlier PZ12,
PZX and OFM results \cite{PeriodZeroOriginal}; here they are applied
at the finite zeros whose existence is newly established by ANP2.
The complete factor test does not need admissibility. Retain the nine
distinct original frequencies and full coefficients in OFM2--3,
\[
 \lambda_{uw}=1+(2u-2)\delta+i(2w-2)\gamma,
 \quad
 f_j(\xi,z_0)=\sum_{u,w=0}^2c_{j,uw}(z_0)e^{\lambda_{uw}\xi}.
 \tag{GFC13}
\]
The nine-by-nine matrix $(\lambda_{uw}^m)_{0\le m\le8}$
is invertible because all frequencies are distinct.
Thus exactly one of the following alternatives holds for each factor:
its first nonzero derivative occurs at some $n_j\in\{0,\ldots,8\}$,
or all nine derivatives vanish and the factor vanishes identically
in $\xi$. In the latter case the invertible derivative matrix forces
every coefficient in GFC13 to vanish. The test is therefore finite
and complete even outside the five-orbit locus.

If at least one factor is identically zero, $E_A(\xi,z_0)$ is
identically zero, all its moments vanish, and
\[
       g(t,z_0)=0,\qquad B_{D,s}(z_0)=0
                       \quad\text{for every }D.
 \tag{GFC14}
\]
Conversely a product of finitely many nonzero entire functions of
$\xi$ cannot vanish identically, by the same isolated-zero argument
used for GFC1. This proves that GFC14's symbol condition occurs
precisely when one of the original factors is identically zero.
It makes no assertion that $\mu_0(z)$ vanishes as a function of
the period: GFC3 has already proved its finite multiplicity there.

If all four factors are nonzero functions of $\xi$, define their
orders by the finite test above. Then at the selected period
\[
 v=n_1+n_2+n_3+n_4\in\{1,\ldots,32\},\qquad
 \mu_v(z_0)=v!\prod_{j=1}^4
                 \frac{\partial_\xi^{n_j}f_j(0,z_0)}{n_j!}\ne0,
 \quad
 g(t,z_0)=t^v\widetilde g(t),\quad
 \widetilde g(0)=\frac{(-i)^v\mu_v(z_0)}{v!}\ne0.
 \tag{GFC15}
\]
The lower bound $v\ge1$ follows from $\mu_0(z_0)=0$.
Multiplication of the four Taylor series proves the order and
coefficient, and $\omega(t)=-it+O(t^3)$ proves the displayed
coefficient of $g$.

Let $J_D$ be the upper shift with ones in positions $(i,i+1)$.
The full finite matrix at this point is
\[
 B_{D,s}(z_0)=R_D^{-1}J_D^vT_D(\widetilde g)R_D.
 \tag{GFC16}
\]
The Toeplitz factor is upper triangular with nonzero diagonal and
is invertible. It commutes with $J_D$ since both are polynomials
in that shift. Therefore, putting $v_D=\min(v,N)$,
\[
 \begin{split}
 \operatorname{rank}B_{D,s}(z_0)&=N-v_D,\\
 \ker B_{D,s}(z_0)&=\mathcal P_{v_D-1}
               \quad\text{inside its original source frame},\\
 \operatorname{im}B_{D,s}(z_0)&=\mathcal P_{D-v_D}
               \quad\text{inside its original target frame}.
 \end{split}\tag{GFC17}
\]
Here $\mathcal P_{-1}=0$. The kernel statement follows either
from the shift factor or from the exact translation formula
$T_AS^m=\sum_{n\ge v}\binom mn\mu_nS^{m-n}$; its leading
coefficient is nonzero when $m\ge v$.
If $v>D$, GFC17 correctly says that this cutoff matrix is zero
although the complete symbol is not identically zero.

For finite $v$, the actual order-adjusted map remains the distinct
isomorphism
\[
 (\mathcal P_{D+v}/\mathcal P_{v-1},\|\cdot\|_{s,c_0,\rm quot})
       \longrightarrow(\mathcal P_D,\|\cdot\|_{s,c_1}),\qquad
 \widetilde B_{D,s}=
 \operatorname{diag}(\rho_0,\ldots,\rho_D)^{-1}
 T_D(\widetilde g)
 \operatorname{diag}(\rho_v,\ldots,\rho_{v+D}).
 \tag{GFC18}
\]
Its nonzero triangular diagonal proves invertibility with the complete
quotient norm. When the symbol vanishes identically there is no
finite $v$ or such isomorphism: the original translation annihilates
every polynomial, as all its moments are zero. Thus neither possible
symbol case is replaced by an assumption of admissibility.

The rank formula sharpens GFC12 where its hypothesis is decided by
GFC13: for finite $v\le D$ and $r\ge v$, the nonzero exterior
power $\bigwedge^{N-r}B(z_0)$ gives $m_{N-r}=0$. Hence
\[
 p_r=\nu N,\qquad
 \lim_{z\to z_0}|z-z_0|^{\nu N}\|\bigwedge^rB(z)^{-1}\|
   =\frac{\|\bigwedge^{N-r}B(z_0)\|}{|c|^N}>0
              \quad(v\le r\le N).
 \tag{GFC19}
\]
For $r<v$ the minors vanish at $z_0$, and the finite test GFC10
gives their exact nonzero higher jets rather than presuming simple
rank loss. If $B(z_0)=0$, it gives $m_k\ge k$ for $k\ge1$,
since each $k$-minor is a sum of products of $k$ entries vanishing
at that point. Together with GFC10 this proves
$\nu r\le p_r\le\nu N-(N-r)$ for $1\le r<N$ in that case,
while the top exterior identity is unchanged.

\subsection{Exact original-period constants and the scope of the conclusion}
Put $u_0=1/z_0\ne0$ and use $z=1/u$ on every original logarithm
branch. The factor functions and GFC6 are independent of that branch
by the full diagonal cancellation PCL7--9. The exact coordinate identity is
\[
 z-z_0=-\frac{u-u_0}{u u_0},\qquad
 \mu_0(u^{-1})=
       c\left(-\frac1{u_0^2}\right)^\nu
                  (u-u_0)^\nu+O((u-u_0)^{\nu+1}).
 \tag{GFC20}
\]
Therefore the complete constants in the original period coordinate are
\[
 \boxed{\quad
 \lim_{u\to u_0}|u-u_0|^{p_r}
       \|\bigwedge^rB_{D,s}(u^{-1})^{-1}\|
   =|u_0|^{2p_r}\frac{\|L_{N-r}\|}{|c|^N}>0,
 \quad 1\le r\le N.
 \quad}\tag{GFC21}
\]
For the top exterior rank this is exactly
$|u_0|^{2\nu N}|c|^{-N}$ with exponent $\nu N$.

Consequently every fixed original quartet and original nonzero unit
has at least one finite nonzero complex period at which the old
fixed-space conductor is singular, and every inverse exterior rank
is unbounded when that period is approached through the same
coefficient family. The symbol alternatives GFC13--18 determine
precisely what the complete map does at the selected period.
This obstructs a uniform fixed-space inverse-exterior bound across
all such complex periods. It does not identify those periods with
the arithmetic parameter set, impose the five-orbit condition on
all selected zeros, or change the order-adjusted spaces in GFC18.
